\begin{center}
    {\LARGE \textbf{Abstract}}
\end{center}
\addcontentsline{toc}{chapter}{Abstract}

Cattle farming is one of the fundamental pillars of the Brazilian economy. Bovine carcass ultrasonography makes it possible to measure, \textit{in vivo}, the Ribeye Area (REA), Backfat Thickness (BFT), and Rump Fat Thickness (RFT), which are related to muscularity, early fat development in the carcass, and meat quality. However, the manual delineation of these indicators in ultrasound images is a subjective, time-consuming process that depends on the specialist’s experience. The objective of this work is to develop and compare two models to automatically delimit and measure these indicators in ultrasound images from two cattle breeds. The first model is a U-Net convolutional neural network, and the second, U-Net Fuzzy, is a U-Net integrated with the Adaptive Neuro-Fuzzy Inference System (ANFIS). The images were obtained from live animals at the Advanced Center for Technological Research in Agribusiness (CAPTA) in Sertãozinho-SP, totaling 135 images, of which 108 are Nelore and 27 Caracu. These images were expanded using data augmentation techniques, such as adding versions with reduced and increased brightness. The segmentation masks used as reference, known as \textit{Ground Truth}, were manually created with the assistance of a specialist, consisting of binary images used to isolate specific regions within the image, serving as a reference for the training and evaluation of the models. The performance of the models was evaluated using metrics such as accuracy, Mean Absolute Error (MAE), and the Dice coefficient, which quantifies the similarity between the segmentation produced by the model and the segmentation masks created by the specialist. The results indicate that the U-Net Fuzzy model outperformed the U-Net model in measuring the REA and RFT indicators across all evaluation metrics, especially achieving accuracies of $94.17\%$ for REA and $98.6\%$ for RFT. In contrast, the U-Net model achieved better results for all metrics regarding the BFT indicator, with particular emphasis on the Dice coefficient metric of $70.25\%$. Finally, a computational application was developed capable of detecting specific regions and automatically measuring the indicators of interest in bovine ultrasound images using both models, potentially assisting in herd selection and management processes.
\vspace{1cm}
\\
\textbf{Keywords}: Bovine ultrasonography; Image delineation; Neural networks; U-Net; ANFIS.